\documentclass[../DoAn.tex]{subfiles}
\begin{document}

\section{Môi trường thực nghiệm, bộ dữ liệu và tham số đánh giá}
\label{sec:tn-moitruong}

\subsection{Môi trường phần cứng và phần mềm}
\label{subsec:tn-hardware}

Toàn bộ thực nghiệm được thực hiện trên một máy tính cá nhân với cấu hình phổ thông, không có GPU chuyên dụng --- phản ánh điều kiện phần cứng thực tế mà đề tài hướng tới khả năng triển khai:

\begin{itemize}
    \item \textbf{Phần cứng:} CPU Intel Core i7 thế hệ 11, RAM 16 GB DDR4, ổ SSD. Toàn bộ quá trình huấn luyện chạy trên CPU với PyTorch CPU backend.
    \item \textbf{Hệ điều hành:} Windows 11 (host) với WSL2 Ubuntu 22.04 LTS. WSL2 cấu hình chế độ Mirrored Networking để chia sẻ địa chỉ IP với Windows host, cho phép traffic từ mạng LAN vật lý tiếp cận trực tiếp interface WSL.
    \item \textbf{Thư viện chính:} Python 3.8, PyTorch 2.0 (FT-Transformer), TensorFlow/Keras 2.12 (Autoencoder), LightGBM 4.0, scikit-learn 1.3, imbalanced-learn 0.11.
    \item \textbf{Công cụ mạng:} CICFlowMeter v4 (trích xuất đặc trưng flow từ live interface, timeout 30 giây), Snort 3.1 (IDS rule-based, Custom Rules + Community Rules + Emerging Threats Open).
\end{itemize}

\subsection{Bộ dữ liệu thực nghiệm}
\label{subsec:tn-datasets}

Ba bộ dữ liệu được sử dụng qua từng giai đoạn nghiên cứu, phản ánh quá trình dịch chuyển từ benchmark chuẩn sang dữ liệu thu thập thực chiến:

\begin{table}[H]
    \centering
    \caption{Tổng quan các bộ dữ liệu thực nghiệm}
    \label{tab:datasets-overview}
    \begin{tabular}{llrrr}
        \hline
        \textbf{Dataset} & \textbf{Giai đoạn} & \textbf{Tổng mẫu} & \textbf{Đặc trưng} & \textbf{Số lớp} \\
        \hline
        NSL-KDD Train+  & 1 & 125.973 & 41 gốc $\to$ 122 & 5 \\
        NSL-KDD Test+   & 1 & 22.542  & 41 gốc $\to$ 122 & 5 \\
        CIC-IDS-2017    & 2, 3 & 2.829.385 & 77 & 15 \\
        Testbed Combined & 3 & 111.825 & 80 & 5 \\
        \hline
    \end{tabular}
\end{table}

NSL-KDD sử dụng 41 đặc trưng gốc, được mở rộng lên 122 chiều sau mã hóa one-hot cho ba đặc trưng phân loại (\texttt{protocol\_type}: 3 giá trị, \texttt{service}: 70 giá trị, \texttt{flag}: 11 giá trị) và MinMaxScaler cho phần còn lại. CIC-IDS-2017 sử dụng 77 đặc trưng --- gồm 68 đặc trưng thống kê CICFlowMeter, 5 chỉ báo nhóm cổng và 4 đặc trưng tỉ lệ được kỹ thuật hóa trong pipeline tiền xử lý --- chuẩn hóa bằng PowerTransformer (Yeo-Johnson). Tập Testbed Combined mở rộng tập này lên 80 đặc trưng bằng 3 đặc trưng từ bước Model Surgery (IAT\_CV, Bwd\_Pkt\_Ratio, Pkt\_Size\_Ratio), chuẩn hóa bằng PowerTransformer được fit trên phân phối Testbed.

\subsection{Các chỉ số đánh giá}
\label{subsec:tn-metrics}

Tất cả các bài toán phân loại trong đề tài đều có mất cân bằng lớp đáng kể, khiến accuracy trở thành chỉ số không đáng tin cậy --- một mô hình dự đoán tất cả là Benign cũng đạt trên 80\% accuracy trên CIC-IDS-2017. Ba chỉ số chính được sử dụng:

\textbf{Macro-F1} là metric chính để so sánh mô hình, trong đó mỗi lớp được trọng số bằng nhau, phản ánh yêu cầu bảo mật: bỏ sót một cuộc tấn công U2R (52 mẫu training) nghiêm trọng tương đương bỏ sót DoS (45.927 mẫu training):
\begin{equation}
    \text{Macro-F1} = \frac{1}{K}\sum_{c=1}^{K}\frac{2 P_c R_c}{P_c + R_c}
    \label{eq:macro-f1}
\end{equation}

\textbf{MCC} (Matthews Correlation Coefficient) được dùng cho đánh giá domain adaptation vì nhạy hơn với hiệu năng đồng thời trên tất cả các lớp khi phân phối bất đối xứng cực đoan. Giá trị MCC = 0 tương đương dự đoán ngẫu nhiên; MCC âm cho thấy mô hình đang phân loại sai hướng có hệ thống.

\textbf{Balanced Accuracy} (= trung bình Recall qua tất cả các lớp) được dùng bổ sung cho Testbed để phản ánh rõ hiệu năng trên từng lớp thiểu số.

\section{Kết quả giai đoạn NSL-KDD}
\label{sec:tn-nslkdd}

Giai đoạn NSL-KDD thiết lập hai thành phần --- Autoencoder Anomaly Gate và Stacking Ensemble --- đồng thời xác định các chiến lược xử lý mất cân bằng cốt lõi. Tất cả kết quả trong mục này được đánh giá trên KDDTest+ (22.542 mẫu) theo Macro-F1 trên 5 lớp.

Một lưu ý về phạm vi đo, cần thiết để đọc đúng các bảng phía dưới: toàn bộ số liệu ở đây đến từ \textbf{cấu hình Stacking phẳng 5 lớp}, trong đó mô hình dự đoán trực tiếp cả nhãn Normal lẫn bốn nhãn tấn công trên toàn bộ KDDTest+. Cấu hình này \textit{không} đi qua Autoencoder Gate. Gate là thành phần của bản triển khai inference (nối tiếp một bộ phân loại 4 lớp tấn công) và được đánh giá riêng bằng chỉ số nhị phân. Việc tách bạch này là cần thiết vì hai cấu hình có không gian nhãn khác nhau --- 5 lớp so với 4 lớp --- nên các chỉ số của chúng không thể đặt chung một bảng.

\subsection*{So sánh các phiên bản kiến trúc}

Bảng~\ref{tab:nslkdd-ablation} trình bày quá trình cải thiện kiến trúc từ mô hình đơn lẻ đến Stacking Ensemble cuối:

\begin{table}[H]
    \centering
    \caption{So sánh Macro-F1 trên 5 lớp (KDDTest+) qua các phiên bản kiến trúc, đo ở cấu hình Stacking phẳng.}
    \label{tab:nslkdd-ablation}
    \resizebox{\textwidth}{!}{%
    \begin{tabular}{lcc}
        \hline
        \textbf{Kiến trúc} & \textbf{Macro F1} & \textbf{Ghi chú} \\
        \hline
        FT-Transformer đơn lẻ (baseline) & 0,638 & Không tối ưu val set \\
        FT-Transformer + Boost-Minority Val & 0,654 & Best FT-Transformer \\
        LightGBM đơn lẻ & 0,668 & Probe/U2R tốt hơn; R2L kém hơn \\
        \textbf{Stacking Ensemble (FTT + LGBM → Meta-LR)} & \textbf{0,681} & \textbf{Kết quả cuối} \\
        \hline
    \end{tabular}%
    }
\end{table}

\begin{figure}[H]
    \centering
    \includegraphics[width=0.72\textwidth]{fig_nslkdd_ablation}
    \caption{So sánh Macro F1 qua các phiên bản kiến trúc trên NSL-KDD (5 lớp, KDDTest+).}
    \label{fig:nslkdd-ablation}
\end{figure}

Mỗi bước cải thiện có động cơ rõ ràng: Boost-Minority Val giải quyết vấn đề checkpoint selection trên U2R (chỉ 52 training samples); LightGBM khai thác phân vùng đặc trưng cứng hiệu quả cho Probe; Stacking Ensemble kết hợp điểm mạnh của cả hai --- FT-Transformer xử lý tốt R2L (Recall cao hơn nhờ Attention cross-feature), LightGBM xử lý tốt Probe và U2R.

\subsection*{Kết quả chi tiết per-class của kiến trúc cuối}

\begin{table}[H]
    \centering
    \caption{Kết quả phân loại 5 lớp trên KDDTest+ --- Stacking Ensemble (FTT + LightGBM $\to$ Meta-LR), không qua Autoencoder Gate.}
    \label{tab:nslkdd-final}
    \begin{tabular}{lcccc}
        \hline
        \textbf{Lớp} & \textbf{Precision} & \textbf{Recall} & \textbf{F1} & \textbf{Support} \\
        \hline
        Normal  & 0,7154 & 0,9682 & 0,8228 & 9.711  \\
        DoS     & 0,9630 & 0,8032 & 0,8759 & 7.458  \\
        Probe   & 0,8024 & 0,7852 & 0,7937 & 2.421  \\
        R2L     & 0,9681 & 0,2523 & 0,4003 & 2.885  \\
        U2R     & 0,5517 & 0,4776 & 0,5120 & 67     \\
        \hline
        \textbf{Macro avg} & \textbf{0,8001} & \textbf{0,6573} & \textbf{0,6809} & 22.542 \\
        Accuracy        & &  & 0,8009 & 22.542 \\
        \hline
    \end{tabular}
\end{table}

\begin{figure}[H]
    \centering
    \includegraphics[width=0.85\textwidth]{fig_nslkdd_perclass}
    \caption{Kết quả phân loại 5 lớp trên KDDTest+ --- Stacking Ensemble (FTT + LightGBM $\to$ Meta-LR).}
    \label{fig:nslkdd-perclass}
\end{figure}

\subsection*{Phân tích kết quả}

\textbf{Điểm mạnh:} Normal Recall đạt 0,9682 --- mô hình nhận ra gần như trọn vẹn lưu lượng bình thường, vốn chiếm 9.711 trong 22.542 mẫu của KDDTest+. Cần đọc con số này cùng với Precision tương ứng chỉ 0,7154: mô hình gán nhãn Normal khá rộng tay, nên bắt được hầu hết Normal thật nhưng cũng kéo theo một lượng đáng kể mẫu tấn công bị gán nhầm thành Normal --- chủ yếu là R2L, như phân tích ngay dưới đây. DoS F1 = 0,8759 và Probe F1 = 0,7937 phản ánh các lớp này có pattern flow đặc trưng và đủ dữ liệu huấn luyện. U2R F1 = 0,5120 với chỉ 52 mẫu huấn luyện là kết quả cạnh tranh so với các công bố trên cùng bộ dữ liệu.

\textbf{Điểm yếu cấu trúc --- R2L:} R2L Recall = 0,2523 là giới hạn khó vượt qua, và điều đáng chú ý là nó đi kèm Precision rất cao (0,9681). Cặp chỉ số này mô tả một mô hình \textit{quá thận trọng} với R2L: khi đã gọi một flow là R2L thì gần như luôn đúng, nhưng ba phần tư số R2L thật bị bỏ sót và rơi vào nhãn Normal.

Nguyên nhân không nằm ở thuật toán mà ở chính cấu tạo của bộ dữ liệu. R2L là các tấn công truy cập trái phép từ xa, biểu hiện ở mức flow gần như trùng khớp với một phiên đăng nhập hợp lệ --- cùng giao thức, cùng khoảng kích thước gói, cùng nhịp thời gian --- nên phần lớn tín hiệu phân biệt nằm ở nội dung payload chứ không ở thống kê flow. Nghiêm trọng hơn, tập huấn luyện và tập kiểm tra của NSL-KDD lệch nhau về giao thức: R2L trong KDDTrain+ tập trung ở telnet, trong khi phần lớn R2L của KDDTest+ đi qua pop3. Mô hình vì vậy được yêu cầu tổng quát hóa sang một giao thức chưa từng quan sát. Đây là giới hạn cứng của bộ dữ liệu, và cũng là một trong những lý do dẫn tới quyết định chuyển sang CIC-IDS-2017 ở giai đoạn sau.

\textbf{Đặt trong bối cảnh literature:} Macro-F1 = 0,6809 trên KDDTest+ (5 lớp, 5-way end-to-end, không rò rỉ train/test) nằm trong khoảng cạnh tranh 0,60--0,70 của các nghiên cứu đánh giá trung thực. Các kết quả báo cáo 99\%+ trên NSL-KDD trong literature thường sử dụng cross-validation trên KDDTrain+ --- không phản ánh distribution shift thật giữa KDDTrain+ và KDDTest+.

Kết quả này xác nhận hai điều: kiến trúc Two-Stage với Stacking Ensemble là hướng đúng cho bài toán mất cân bằng nhiều lớp; và giới hạn trên của NSL-KDD (dữ liệu 1998, thiếu diversity protocol) không thể vượt qua bằng cải tiến thuật toán đơn thuần --- động cơ trực tiếp để chuyển sang CIC-IDS-2017 với dữ liệu hiện đại và phong phú hơn.

\section{Kết quả Two-Stage Cascade NIDS trên CIC-IDS-2017}
\label{sec:tn-cic}

\subsection{Hiệu quả kiến trúc Two-Stage so với mô hình 1-Stage}
\label{subsec:tn-cic-twostage}

Vấn đề cốt lõi với mô hình 1-Stage trực tiếp trên CIC-IDS-2017: Benign chiếm 80,32\% trong 2,83 triệu flow. Khi một FT-Transformer đơn lẻ phải phân loại 9 lớp đồng thời (Benign và 8 nhóm tấn công gộp từ 15 nhãn gốc), gradient bị dominated bởi Benign và các lớp lớn, khiến các lớp tấn công hiếm như Infiltration (0,0013\%) và Botnet (0,070\%) nhận gradient quá nhỏ để học được pattern phân biệt.

Kiến trúc Two-Stage tách bài toán thành hai không gian riêng biệt:

\textbf{Tầng 1 --- Gating Network (6 lớp):} Phân loại flow thành \{Benign, DoS, DDoS, PortScan, Brute Force, Suspicious\} trên 77 đặc trưng. Bốn lớp tấn công phổ biến được phân loại trực tiếp. Ngưỡng $P(\text{Suspicious}) \geq 85\%$ quyết định flow nào được chuyển tiếp; flow bị đánh dấu Suspicious nhưng dưới ngưỡng được ép về Benign. Như đã phân tích ở Mục~\ref{subsec:gd2-stage1}, ngưỡng này hoạt động như một van giảm false positive --- nó hy sinh một phần recall để đổi lấy precision, chứ không mở rộng vùng nghi ngờ.

\textbf{Tầng 2 --- Expert Network (5 lớp):} Chỉ nhận Suspicious flows từ Tầng 1 và phân loại thành \{Benign, Web Attack, Bot, Infiltration, Heartbleed\}. Khi các lớp đa số (Benign, DoS, DDoS, PortScan, Brute Force) đã được lọc ra ngoài, Expert Network học trong không gian sạch hơn --- gradient tập trung phân biệt \textit{giữa các tấn công phức tạp} thay vì bị overwhelmed bởi lưu lượng đa số.

Cơ chế mang lại lợi ích của phân tầng có thể phát biểu chính xác qua tỉ lệ đóng góp gradient. Trong cấu hình một tầng, mỗi bước cập nhật trọng số nhận tín hiệu từ một batch mà Benign chiếm 80,32\%; các lớp như Infiltration (0,0013\%) và Botnet (0,070\%) đóng góp một phần gradient nhỏ hơn từ ba tới bốn bậc độ lớn, và phần đóng góp đó tiếp tục bị pha loãng bởi bốn lớp tấn công phổ biến vốn có hàng trăm nghìn mẫu. Hệ quả là mô hình hội tụ về một nghiệm tối ưu cho các lớp đa số, trong đó việc luôn dự đoán ``không phải Infiltration'' gần như không bị phạt.

Kiến trúc phân tầng loại bỏ chính sự pha loãng đó. Sau khi Tầng 1 đã tách trực tiếp bốn lớp tấn công phổ biến cùng phần lớn Benign, Expert Network chỉ còn phải phân biệt \textit{trong nội bộ nhóm lớp hiếm}. Tỉ lệ tương đối giữa các lớp mà nó nhìn thấy thay đổi hoàn toàn: một lớp chiếm 0,0013\% toàn tập có thể chiếm vài phần trăm trong luồng Suspicious. Đây là lý do các lớp hiếm học được ở Tầng 2 những pattern mà cấu hình một tầng không học được --- và cũng là lý do Expert Network có thể dùng kiến trúc nhỏ hơn ($d=64$, 3 tầng) mà vẫn hiệu quả, vì nhiệm vụ của nó hẹp hơn hẳn.

\subsection{Đóng góp của Hard Negative Mining}
\label{subsec:tn-cic-hnm}

Sau vòng huấn luyện đầu tiên của Expert Network, Botnet và Infiltration vẫn là hai lớp có hiệu năng thấp nhất. Phân tích confusion matrix cho thấy hai nguyên nhân khác nhau:

\begin{itemize}
    \item \textbf{Botnet}: Nhiều false positive --- nhầm một phần Benign long-idle connection thành Botnet C2 heartbeat do cả hai đều có pattern periodic, low-rate ở tần suất tương đương.
    \item \textbf{Infiltration}: Nhiều false negative --- bỏ sót Infiltration bị phân loại sai thành Benign SSH/HTTP vì flow Infiltration riêng lẻ không có dấu hiệu bất thường rõ ràng ở mức per-flow.
\end{itemize}

Hard Negative Mining (HNM) chạy hai vòng: thu thập các mẫu bị phân loại sai trên tập validation, thêm vào training với trọng số tăng cường $w_\text{hard} = 2{,}0$, huấn luyện lại. Bảng~\ref{tab:hnm-progress} cho thấy kết quả tích lũy:

Để tách bạch đóng góp của HNM khỏi các kỹ thuật xử lý mất cân bằng khác, thực nghiệm so sánh trực tiếp nó với phương án thông dụng nhất --- nhân trọng số lớp --- trên cùng lớp Botnet và cùng điều kiện đo nội bộ Tầng 2.

\begin{table}[H]
    \centering
    \caption{So sánh Hard Negative Mining với nhân trọng số lớp, đo trên Botnet ở phạm vi nội bộ Tầng 2.}
    \label{tab:hnm-progress}
    \begin{tabular}{lc}
        \hline
        \textbf{Chiến lược xử lý mất cân bằng} & \textbf{Botnet F1} \\
        \hline
        Nhân trọng số lớp $\times 10$      & 0,5103 \\
        Hard Negative Mining, 2 vòng       & \textbf{0,6512} \\
        \hline
        Chênh lệch & $+0{,}1409$ \\
        \hline
    \end{tabular}
\end{table}

\begin{figure}[H]
    \centering
    \includegraphics[width=0.65\textwidth]{fig_hnm_ablation}
    \caption{So sánh hiệu quả của Hard Negative Mining và nhân trọng số lớp trên lớp Botnet.}
    \label{fig:hnm-ablation}
\end{figure}

Chênh lệch khoảng 14 điểm tuyệt đối giữa hai chiến lược cho thấy điều quan trọng không nằm ở việc \textit{tăng bao nhiêu} tín hiệu học cho lớp thiểu số, mà ở việc \textit{tăng vào đâu}. Nhân trọng số lớp phân bổ đều mức tăng cường lên mọi mẫu Botnet, trong đó phần lớn là các mẫu mô hình vốn đã phân loại đúng --- phần tăng cường dành cho chúng gần như không mang lại thông tin mới, thậm chí có thể đẩy biên quyết định lấn sang vùng Benign và làm tăng false positive. Hard Negative Mining chỉ tăng cường các mẫu nằm sát biên, tức đúng vùng mà mô hình còn chưa quyết định được.

Cần lưu ý về phạm vi đo: các con số trong Bảng~\ref{tab:hnm-progress} được tính trong nội bộ Tầng 2, trên tập validation riêng của tầng này. Khi đánh giá đầu-cuối trên toàn bộ cascade --- nơi Botnet còn phải đi qua bộ lọc của Tầng 1 --- F1 của lớp này ở cùng phiên bản đạt 0,6294 (Bảng~\ref{tab:res-cic-ablation}). Hai con số đo hai thứ khác nhau và không nên đặt cạnh nhau như một chuỗi cải tiến.

\subsection{Kết quả Ensemble Voting cuối cùng}
\label{subsec:tn-cic-ensemble}

Sau HNM, Expert Network đơn lẻ vẫn còn hai điểm yếu còn dư: Botnet có false positive cao (precision thấp gây alert fatigue), Infiltration vẫn có false negative cao (recall thấp bỏ sót tấn công). Asymmetric Ensemble Voting kết hợp Expert Network (FT-Transformer) với Random Forest (150 cây, max\_depth=25) và KNN (K=16) theo hai luật bỏ phiếu bất đối xứng:

\begin{itemize}
    \item \textbf{Quy tắc Ưu tiên Infiltration:} RF \textit{hoặc} KNN bỏ phiếu Infiltration $\Rightarrow$ kết quả là Infiltration. Vì chi phí bỏ sót Infiltration (hệ thống bị xâm nhập không được phát hiện) lớn hơn nhiều so với false alarm, quy tắc này chấp nhận đánh đổi precision lấy recall.
    \item \textbf{Quy tắc Đồng thuận Botnet:} Kết quả là Botnet \textit{chỉ khi} FT-Transformer dự đoán Botnet và cả RF lẫn KNN không dự đoán Benign. RF và KNN có inductive bias khác FT-Transformer, ít bị nhầm Benign thành Botnet hơn; nếu một trong hai dự đoán Benign thì kết quả bị override thành Benign, giảm mạnh false positive.
\end{itemize}

\begin{table}[H]
    \centering
    \caption{So sánh trước và sau Asymmetric Ensemble Voting, đo đầu-cuối trên 2.529.391 flow của CIC-IDS-2017.}
    \label{tab:cic-final-compare}
    \begin{tabular}{lcc}
        \hline
        \textbf{Chỉ số} & \textbf{Trước Ensemble} & \textbf{Sau Ensemble} \\
        \hline
        Accuracy             & 99,51\%  & \textbf{99,55\%} \\
        Macro F1             & 0,9160   & \textbf{0,9294} \\
        Botnet F1            & 0,6294   & \textbf{0,7344} \\
        \quad Botnet Precision & 47,87\%  & \textbf{79,47\%} \\
        \quad Botnet Recall    & 91,87\%  & 68,26\% \\
        Infiltration F1      & 0,7407   & 0,7407 \\
        \hline
    \end{tabular}
\end{table}

Bảng~\ref{tab:cic-final-compare} cho thấy hai luật bất đối xứng có hiệu lực rất khác nhau trong thực tế, và điều này cần được báo cáo đúng.

Quy tắc Đồng thuận Botnet hoạt động đúng như thiết kế và là nguồn gần như duy nhất của mức tăng Macro F1. Precision của Botnet tăng từ 47,87\% lên 79,47\%: trước đó, cứ hai cảnh báo Botnet thì có hơn một cảnh báo sai. Cái giá phải trả là Recall giảm từ 91,87\% xuống 68,26\% --- khoảng một phần tư số Botnet thật bị hạ xuống Benign vì Random Forest hoặc KNN phản đối. Với một lớp mà false positive gây alert fatigue nghiêm trọng, đánh đổi này là có chủ đích, nhưng nó là đánh đổi thật chứ không phải cải thiện thuần túy.

Quy tắc Ưu tiên Infiltration thì \textbf{không tạo ra thay đổi đo được}: F1 giữ nguyên 0,7407 trước và sau khi áp dụng ensemble. Nguyên nhân là Random Forest và KNN hầu như không đưa ra dự đoán Infiltration nào nằm ngoài tập mà FT-Transformer đã bắt được, nên nhánh phủ quyết gần như không bao giờ được kích hoạt trên tập kiểm định này. Với 33 mẫu Infiltration, tập kiểm định cũng quá nhỏ để kết luận chắc chắn về hiệu lực của quy tắc. Cách diễn giải trung thực là: đây là một biện pháp phòng vệ hợp lý về mặt thiết kế chi phí lỗi, nhưng chưa có bằng chứng thực nghiệm cho thấy nó đóng góp vào kết quả cuối.

\begin{figure}[H]
    \centering
    \includegraphics[width=0.65\textwidth]{fig_ensemble_compare}
    \caption{So sánh Expert Network (sau HNM) và Asymmetric Ensemble cuối cho hai lớp khó nhất.}
    \label{fig:ensemble-compare}
\end{figure}

Cải thiện lớn nhất đến từ đúng hai lớp mà luật bỏ phiếu bất đối xứng nhắm tới: Botnet tăng thêm 0,0832 F1 (Đồng thuận giảm FP) và Infiltration tăng thêm 0,1504 F1 (Ưu tiên tăng Recall). Các lớp còn lại (DoS, DDoS, PortScan, WebAttack) duy trì hiệu năng cao ổn định.

\section{Kết quả thích nghi miền}
\label{sec:tn-domain}

Giai đoạn 3 đánh giá khả năng tái sử dụng mô hình CIC-IDS-2017 trên Testbed WSL --- hai môi trường có hạ tầng mạng, card mạng và hành vi NAT khác nhau. Thực nghiệm được thiết kế theo thứ tự can thiệp tăng dần để định lượng đóng góp riêng lẻ của từng kỹ thuật.

\textbf{Đường cơ sở --- Áp dụng thẳng mô hình CIC lên Testbed:}

Mô hình CIC (Accuracy 99,55\% trên tập kiểm tra CIC) được áp dụng trực tiếp lên 4.542 flow Testbed mà không thay đổi gì. Kết quả: Accuracy = 21,93\%, MCC = $-0{,}015$. MCC âm cho thấy mô hình không chỉ hoạt động kém mà đang phân loại sai hướng có hệ thống --- tệ hơn cả dự đoán ngẫu nhiên. Covariate shift giữa hai môi trường (switch vật lý Gigabit vs card ảo Hyper-V + NAT Windows) đủ lớn để làm mô hình hoàn toàn mất phương hướng.

\textbf{Can thiệp 1 --- Re-fit Scaler đơn thuần (thất bại):}

Giải pháp đơn giản nhất được thử: giữ nguyên trọng số mô hình, chỉ fit lại PowerTransformer trên phân phối Testbed. Kết quả: Accuracy giảm tiếp từ 21,93\% xuống 6,87\%. Can thiệp này tạo mâu thuẫn: mô hình đã học embedding nội bộ dựa trên khoảng giá trị của CIC-scaler cũ, nhưng đầu vào giờ đây có phân phối khác hoàn toàn --- mô hình ``nghe'' ngôn ngữ mới trong khi vẫn ``nghĩ'' bằng ngôn ngữ cũ, tệ hơn ban đầu.

\textbf{Can thiệp 2 --- Layer Freezing Fine-tuning:}

Fine-tuning đồng thời scaler mới lẫn một phần tham số mô hình: đóng băng lớp Feature Embedding và 2 khối Transformer đầu (giữ đặc trưng tổng quát đã học từ CIC), chỉ cập nhật 2 khối Transformer cuối và classification head trên Testbed data. Mixed-Domain Batch (70\% Testbed + 30\% CIC mỗi batch) ngăn Catastrophic Forgetting (CF):
\begin{equation}
    CF = \frac{F1_{\text{CIC,trước}} - F1_{\text{CIC,sau}}}{F1_{\text{CIC,trước}}} \times 100\%
    \label{eq:cf-def}
\end{equation}

\textbf{Can thiệp 3 --- Model Surgery (77$\to$80 đặc trưng):}

Sau Layer Freezing, ba đặc trưng kỹ thuật mới được thêm vào (IAT\_CV, Bwd\_Pkt\_Ratio, Pkt\_Size\_Ratio). Ma trận embedding Feature Tokenizer được mở rộng $\mathbb{R}^{77 \times d} \to \mathbb{R}^{80 \times d}$ với 3 hàng mới khởi tạo ngẫu nhiên, sau đó fine-tuning lại với Layer Freezing.

\begin{table}[H]
    \centering
    \caption{Kết quả qua chuỗi can thiệp domain adaptation. Đánh giá nhị phân (Benign/tấn công) trên 4.542 flow Testbed, trong đó chỉ có 7 flow Benign --- xem lưu ý về phạm vi đo ngay sau bảng.}
    \label{tab:domain-adapt-results}
    \begin{tabular}{lccc}
        \hline
        \textbf{Can thiệp} & \textbf{Accuracy (Testbed)} & \textbf{MCC (Testbed)} & \textbf{CF (\%)} \\
        \hline
        Mô hình CIC áp dụng thẳng  & 21,93\% & $-0{,}015$ & --- \\
        Re-fit Scaler đơn thuần    &  6,87\% & $+0{,}006$ & --- \\
        Layer Freezing Fine-tuning & 99,82\% & $0{,}6825$  & 0,61\% \\
        + Model Surgery (77$\to$80) & 99,87\% & $\mathbf{0{,}7333}$ & 0,61\% \\
        \hline
    \end{tabular}
\end{table}

\begin{figure}[H]
    \centering
    \includegraphics[width=0.92\textwidth]{fig_da_comparison}
    \caption{So sánh Accuracy và MCC qua chuỗi can thiệp domain adaptation (CIC-IDS-2017 → Testbed WSL2).}
    \label{fig:da-comparison}
\end{figure}

Trước khi diễn giải các con số này, cần nêu rõ phạm vi đo, vì nó quyết định mức độ tin cậy có thể gán cho chúng. Chuỗi thí nghiệm domain adaptation được thực hiện ở dạng \textbf{phân loại nhị phân} (Benign so với tấn công), không phải phân loại nhiều lớp như các mục trước. Tập đánh giá gồm \textbf{4.542 flow Testbed, trong đó chỉ có 7 flow Benign}. Sự chênh lệch cực đoan này là hệ quả của cách thu dữ liệu --- các phiên thu tập trung vào lưu lượng tấn công --- và nó đặt ra hai giới hạn phải tính đến khi đọc bảng.

Thứ nhất, chỉ số Accuracy gần như vô nghĩa ở đây: một mô hình dự đoán mọi flow là tấn công cũng đạt trên 99,8\%. Đây chính là lý do MCC được chọn làm chỉ số chính cho mục này. Thứ hai, ngay cả MCC cũng rất nhạy với vài mẫu: khoảng cách giữa Layer Freezing (0,6825) và Model Surgery (0,7333) tương ứng với việc \textit{hai} flow được phân loại lại đúng trong tổng số 4.542 --- ma trận nhầm lẫn chuyển từ 8 xuống 6 lỗi. Chênh lệch đó nằm trong biên độ dao động thống kê của một tập nhỏ như vậy và không nên được trình bày như một cải thiện đã được khẳng định.

Với những giới hạn đó, kết luận đáng tin cậy của chuỗi thí nghiệm này là ở \textit{dấu hiệu} chứ không phải ở \textit{độ lớn}. Khoảng cách giữa nhóm can thiệp thất bại (MCC âm hoặc gần 0) và nhóm can thiệp thành công (MCC trên 0,68) lớn tới mức không thể quy cho nhiễu thống kê: Layer Freezing thực sự khôi phục được năng lực phân biệt mà chuyển miền trực tiếp đã đánh mất. Ngược lại, Re-fit Scaler đơn thuần cho MCC $\approx +0{,}006$ --- gần như không phân biệt được gì --- và kém hơn cả mô hình chuyển thẳng, xác nhận mâu thuẫn giữa scaler và embedding đã phân tích ở Mục~\ref{subsec:gd3-scaler}. Riêng đóng góp bổ sung của Model Surgery thì cần một tập đánh giá lớn hơn và cân bằng hơn mới kết luận được; bằng chứng thuyết phục hơn cho ba đặc trưng mới nằm ở kết quả nhiều lớp của mô hình cuối (Mục~\ref{sec:tn-testbed}), nơi tập kiểm định có 11.183 flow với đủ năm lớp.

Chuỗi can thiệp này cung cấp một quy trình thực tiễn rõ ràng cho việc tái sử dụng mô hình IDS qua các môi trường mạng khác nhau: (1) không thể thay đổi scaler độc lập với tham số mô hình; (2) Layer Freezing + Mixed-Domain Batch là kỹ thuật hiệu quả cho domain adaptation IDS; (3) thiết kế đặc trưng nhắm vào đặc điểm hạ tầng của môi trường đích có thể bổ sung cải thiện đáng kể.

\section{Kết quả mô hình trên Testbed nội bộ}
\label{sec:tn-testbed}

\subsection{Bằng chứng định lượng về tính không phân tách của PortScan và Benign}
\label{subsec:tn-testbed-portscan}

Kết quả đầu tiên khi huấn luyện mô hình trên dữ liệu Testbed nội bộ: PortScan F1 = 7,5\%. Phản ứng ban đầu là điều chỉnh thuật toán: tăng class weight, Focal Loss với $\gamma$ cao hơn, SMOTE agressive. Không có can thiệp nào cải thiện đáng kể. Bảng~\ref{tab:portscan-algo} trình bày kết quả thử nghiệm có kiểm soát:

\begin{table}[H]
    \centering
    \caption{PortScan F1 trên dữ liệu Testbed WSL với các thuật toán và cấu hình khác nhau}
    \label{tab:portscan-algo}
    \begin{tabular}{lcc}
        \hline
        \textbf{Thuật toán} & \textbf{Cấu hình} & \textbf{PortScan F1} \\
        \hline
        FT-Transformer & CE Loss, class\_weight=1.0 & 7,5\% \\
        FT-Transformer & Focal Loss $\gamma = 2$   & 6,8\% \\
        FT-Transformer & class\_weight PortScan $\times$5 & 8,1\% \\
        Random Forest  & 150 cây, default         & 5,2\% \\
        KNN            & K=5                      & 4,9\% \\
        \hline
        \textit{Tất cả} & CIC PortScan inject (5.000 flow) & \textbf{99,9\%} \\
        \hline
    \end{tabular}
\end{table}

\begin{figure}[H]
    \centering
    \includegraphics[width=0.82\textwidth]{fig_portscan_inseparable}
    \caption{PortScan F1 trên môi trường WSL/NAT với các thuật toán khác nhau — kết quả dưới 9\% nhất quán trên mọi mô hình cho thấy vấn đề nằm ở dữ liệu, không phải thuật toán.}
    \label{fig:portscan-inseparable}
\end{figure}

Kết quả $< 9\%$ nhất quán trên mọi thuật toán và mọi cấu hình là bằng chứng mạnh: vấn đề không nằm ở mô hình mà ở dữ liệu. Trong không gian 80 chiều CICFlowMeter, PortScan flow trên WSL/NAT và Benign HTTP flow không phân biệt được với nhau.

Nguyên nhân gốc rễ: nmap gửi SYN packet đến từng port, NAT Windows xử lý mỗi connection riêng lẻ và CICFlowMeter chỉ capture $\approx 7$ flow mỗi lần scan --- mỗi flow là TCP SYN ngắn không có payload, flow statistics (duration $\approx 0$, bytes $\approx 0$, IAT ngẫu nhiên) hoàn toàn trùng với Benign health check hay short HTTP connection.

Sau khi inject 5.000 PortScan flow từ CIC-IDS-2017 Friday (thu thập trên switch vật lý, có nhiều SYN flags, duration đồng nhất, packet count cao): PortScan F1 tăng từ 7,5\% lên \textbf{99,9\%} và duy trì ổn định. Kết quả này xác nhận: khi môi trường thu thập tạo ra dữ liệu không phân tách, giải pháp phải ở tầng dữ liệu --- dùng surrogate data từ môi trường phần cứng thực --- không phải tầng thuật toán.

\subsection{Kết quả phiên bản cuối}
\label{subsec:tn-testbed-final}

Mô hình FT-IDS được đánh giá trên tập validation 11.183 flows với thành phần đa dạng miền: BruteForce validation bao gồm cả hydra SSH/FTP (Testbed nội bộ) và CIC Patator --- đảm bảo đánh giá khả năng tổng quát hóa thực sự, không phải in-distribution.

\begin{table}[H]
    \centering
    \caption{Kết quả FT-IDS trên tập validation đa dạng miền}
    \label{tab:v85-final}
    \resizebox{\textwidth}{!}{%
    \begin{tabular}{lcccc}
        \hline
        \textbf{Lớp} & \textbf{Precision} & \textbf{Recall} & \textbf{F1} & \textbf{Nguồn validation} \\
        \hline
        Benign     & 0,87 & 0,95 & 0,91 & Testbed nội bộ \\
        BruteForce & 0,83 & 0,90 & \textbf{0,86} & hydra + CIC Patator \\
        DoS        & 0,98 & 0,89 & 0,93 & Testbed nội bộ \\
        PortScan   & 1,00 & 1,00 & \textbf{1,00} & CIC Friday (surrogate) \\
        Web Attack & 0,96 & 0,81 & 0,88 & Testbed nội bộ \\
        \hline
        \textbf{Macro F1} & & & \textbf{0,917} & \\
        Balanced Accuracy & & & \textbf{0,9114} & \\
        \hline
    \end{tabular}%
    }
\end{table}

\begin{figure}[H]
    \centering
    \includegraphics[width=0.72\textwidth]{fig_v85_perclass}
    \caption{FT-IDS --- F1-score per-class trên tập validation đa dạng miền (11.183 flows). Đường ngang thể hiện Macro F1 = 0,917.}
    \label{fig:v85-perclass}
\end{figure}

Kết quả đáng chú ý nhất là BruteForce F1 = 86\% (Recall = 90\%) trên tập validation đa dạng --- so với con số 100\% F1 trên tập validation đồng nhất (chỉ CIC Patator) của phiên bản trước. Khoảng cách 14\% tuyệt đối phản ánh chính xác mức độ inflation metric do đánh giá in-distribution, không phải do mô hình thật sự tốt hơn.

Web Attack F1 = 88\% với Recall = 81\%: 19\% flow Web Attack bị bỏ sót chủ yếu là SQLi payload ngắn và XSS thất bại có flow statistics gần giống Benign HTTP POST. Đây là giới hạn của phương pháp flow-based --- payload-level analysis của Snort bổ sung phát hiện chính xác những gì ML bỏ sót.

FT-IDS đạt Macro F1 = 91,66\% trên 5 lớp tấn công trong điều kiện validation đa dạng miền, trên phần cứng phổ thông không có GPU, với dữ liệu thu thập thực từ tấn công xuyên thiết bị trên mạng LAN gia đình.

\section{Thử nghiệm hệ thống End-to-End}
\label{sec:tn-e2e}

Hệ thống Hybrid IDS được đánh giá ở hai mức bổ sung cho nhau. Mức thứ nhất đối chiếu khả năng phát hiện của tầng dấu hiệu (Snort) và tầng học máy (FT-IDS) trên từng loại tấn công, kiểm chứng giả thuyết hai tầng bù điểm mù cho nhau. Mức thứ hai đo hiệu năng phân loại định lượng của hệ thống phát hiện theo phát lại (Mục~\ref{subsec:deploy-replay}) trên toàn bộ corpus Testbed. Đường xử lý thời gian thực hoàn toàn chưa được hiện thực hóa trong phạm vi đồ án; các số liệu của tầng học máy dưới đây đo trên hệ thống replay, dùng đúng mô hình, bộ tiền xử lý và logic quyết định mà phiên bản thời gian thực sẽ sử dụng.

\subsection{Khả năng phát hiện theo loại tấn công}
\label{subsec:e2e-complementarity}

\textbf{Cấu hình:} Năm phiên thu cô lập --- mỗi loại tấn công một phiên riêng, máy tấn công IP \texttt{192.168.0.106}, victim \texttt{192.168.0.103} --- cung cấp corpus phát lại, mỗi tệp lấy tối đa 8.000 flow để cân bằng giữa các lớp. Tầng học máy phân loại các flow này qua mô hình FT-IDS; tầng Snort 3.1 (bộ luật Custom Rules kết hợp Community Rules và Emerging Threats Open) được đối chiếu trên cùng các kịch bản tấn công. Nhãn đối chứng gán theo IP nguồn.

\begin{table}[H]
    \centering
    \caption{Kết quả phát hiện End-to-End theo kịch bản tấn công}
    \label{tab:e2e-results}
    \begin{tabular}{lccp{5.2cm}}
        \hline
        \textbf{Loại tấn công} & \textbf{Snort} & \textbf{ML} & \textbf{Ghi chú} \\
        \hline
        PortScan (nmap SYN)    & \checkmark & \checkmark & Snort: ngay lập tức (per-packet); ML: sau khi flow hoàn chỉnh \\
        BruteForce (hydra SSH) & \checkmark & \checkmark & Snort: rule ssh-brute; ML: qua flow authentication pattern \\
        DoS (slowhttptest)     & $\circ$    & \checkmark & Snort phát hiện một phần; ML phát hiện qua flow duration bất thường \\
        Web Attack --- SQLi    & \checkmark & \checkmark & Cả hai phát hiện; Snort qua HTTP payload, ML qua flow statistics \\
        Web Attack --- XSS     & \checkmark & $\circ$    & Snort phát hiện signature; ML bỏ sót XSS payload ngắn (Recall 81\%) \\
        Benign                 & $\times$   & $\circ$    & Snort: không cảnh báo; ML: một phần Benign SSH/FTP bị nhầm BruteForce (Precision 83\%) \\
        \hline
    \end{tabular}
\end{table}

\textit{Ký hiệu: \checkmark = phát hiện đúng, $\circ$ = phát hiện một phần, $\times$ = không cảnh báo (đúng với lưu lượng bình thường).} Cột ML phản ánh kết quả của FT-IDS trên corpus phát lại; cột Snort phản ánh độ phủ luật trên cùng tập tấn công.

\textbf{Tính bổ sung giữa hai tầng phát hiện:}

Kết quả End-to-End cho thấy Snort và FT-Transformer không cạnh tranh nhau mà điền vào điểm mù của nhau. Snort hoạt động ở tầng gói tin (per-packet, real-time, latency cỡ microsecond), phát hiện ngay khi packet đầu tiên khớp rule. FT-Transformer hoạt động ở tầng flow (per-flow, cần đợi flow hoàn chỉnh hoặc timeout), phát hiện dựa trên pattern thống kê tích lũy.

Trường hợp Web Attack XSS minh họa rõ nhất lý do cần cả hai: XSS payload ngắn tạo ra flow statistics gần giống Benign HTTP POST (tổng bytes tương đương, connection time bình thường) khiến ML bỏ sót; nhưng Snort với rule \texttt{alert http ... content:"\textless script\textgreater"} bắt được ngay ở tầng payload. Ngược lại, DoS slowhttptest tạo ra HTTP connection hoàn toàn hợp lệ về mặt gói tin (không có signature bất thường) nên Snort phát hiện không hoàn toàn, trong khi FT-Transformer phát hiện qua flow duration kéo dài bất thường và IAT đặc trưng của slow-connection attack.

Độ trễ phát hiện của ML: BruteForce và PortScan được phát hiện nhanh (mỗi connection đóng sau vài giây); DoS slowhttptest cần đợi đến khi CICFlowMeter timeout flow (tối đa 30 giây). Đây là sự đánh đổi cố hữu của phương pháp flow-based --- không thể tránh hoàn toàn mà chỉ có thể giảm bằng cách điều chỉnh timeout và kết hợp với Snort real-time.

Đối chiếu hai tầng qua corpus phát lại xác nhận giả thuyết cốt lõi của đề tài: kiến trúc Hybrid IDS (signature-based + ML-based) đạt độ bao phủ tấn công rộng hơn bất kỳ hướng tiếp cận đơn lẻ nào, và sự kết hợp hai tầng này không thể thay thế cho nhau --- mỗi tầng là điều kiện cần để bù đắp điểm mù của tầng còn lại.

\subsection{Hiệu năng định lượng trên corpus phát lại}
\label{subsec:tn-replay-quant}

Trên corpus gộp năm phiên (40.000 flow, 8.000 mỗi loại; PortScan dùng bản thu lại trên Windows host, xem Mục~\ref{subsec:tn-baihoc}), hiệu năng của hệ thống được nâng dần qua hai can thiệp hậu xử lý, đo bằng Macro F1 để phạt đều mọi lớp (Bảng~\ref{tab:replay-progress}).

\begin{table}[H]
    \centering
    \caption{Tiến triển Macro F1 của hệ thống replay theo từng can thiệp hậu xử lý}
    \label{tab:replay-progress}
    \begin{tabular}{lccc}
        \hline
        \textbf{Cấu hình} & \textbf{Macro F1} & \textbf{Benign Recall} & \textbf{PortScan F1} \\
        \hline
        FT-IDS, không hậu xử lý                     & 0,679 & 75,4\% & 0,000 \\
        \quad + ngưỡng tin cậy $T=0{,}6$          & 0,725 & 95,2\% & 0,000 \\
        \quad + luật quét cổng                    & \textbf{0,978} & \textbf{95,2\%} & \textbf{0,996} \\
        \hline
    \end{tabular}
\end{table}

Ở cấu hình cuối, mọi lớp đạt F1 không dưới 0,95: PortScan 0,996, BruteForce 0,991, DoS 0,951, Web Attack 0,990, Benign 0,961; accuracy tổng 0,977.

Macro F1 = 0,978 trên corpus phát lại và Macro F1 = 91,66\% mà FT-IDS đạt trên tập validation đa dạng miền (Mục~\ref{sec:tn-testbed}) không so sánh trực tiếp được: chúng đo trên hai tập dữ liệu khác nhau, và 0,978 còn bao gồm chuỗi luật hậu xử lý mà 91,66\% không có. Trên chính corpus phát lại, mô hình đơn lẻ chỉ đạt Macro F1 = 0,679 (Bảng~\ref{tab:replay-progress}); hai luật hậu xử lý nâng nó lên 0,978 --- ngưỡng tin cậy đưa Benign Recall từ 75,4\% lên 95,2\%, và luật quét cổng đưa PortScan từ không phát hiện được (F1 0,000) lên 0,996.

\subsection{Lỗi sai và bài học}
\label{subsec:tn-baihoc}

Quá trình xây dựng hệ thống replay bộc lộ một số lỗi đặc thù của môi trường ảo hóa và NAT, cùng vài nguyên tắc kỹ thuật rút ra từ việc sửa chúng.

\textbf{Gán nhãn theo IP phải dùng đúng IP sau NAT.} Lưu lượng từ WSL bị NAT qua máy Windows host, nên máy victim quan sát được địa chỉ nguồn là IP LAN của host (\texttt{192.168.0.106}), không phải IP nội bộ của WSL. Lần gán nhãn đầu dùng IP WSL khiến toàn bộ flow tấn công bị gán sai. Nhãn đối chứng phải lấy theo IP nguồn quan sát được tại điểm thu, không phải IP danh nghĩa của máy phát.

\textbf{PortScan phải bắt trên máy Windows host, không phải trong WSL.} Gói tới cổng đóng được chính Windows host trả RST và không bắc cầu vào không gian mạng của WSL; tcpdump trong WSL vì vậy chỉ thấy số ít cổng mở. Mất đi dải cổng bị quét, một phiên PortScan thu trong WSL trông giống DoS --- nhiều flow ngắn tới ít cổng --- và đây chính là nguyên nhân gốc khiến PortScan F1 chỉ đạt 0,000 ở cấu hình đầu. Thu lại trên host khôi phục đặc trưng quét cổng: số cổng đích duy nhất tăng từ 6--9 lên hơn 22.000.

\textbf{nmap cần tùy chọn \texttt{-Pn}.} Khi victim chặn ping, nmap mặc định coi host đã tắt và bỏ qua bước quét, chỉ để lại vài flow rải rác. Bỏ bước kiểm tra sống bằng \texttt{-Pn} là điều kiện để thu được một phiên quét đầy đủ.

\textbf{Một luật dấu hiệu có thể bù đặc trưng còn thiếu rẻ hơn tái huấn luyện.} Bộ 80 đặc trưng của FT-IDS không mã hóa mức độ quét cổng, nên mô hình không thể tách PortScan khỏi DoS. Một luật đếm cổng đích theo nguồn --- đúng tinh thần dấu hiệu --- bù trực tiếp điểm mù đó, nâng PortScan F1 từ 0,000 lên 0,996 mà không chạm tới trọng số mô hình. So với việc thêm đặc trưng thứ 81 rồi huấn luyện lại toàn bộ, luật hậu xử lý đạt cùng mục tiêu với chi phí thấp hơn nhiều và giữ nguyên hành vi đã kiểm chứng của mô hình trên bốn lớp còn lại.

\textbf{Đo trước khi sửa.} Ảnh chụp chỉ số ở mốc gốc cho thấy điểm kéo Macro F1 xuống mạnh nhất là PortScan (F1 0,000), không phải false positive của Benign như giả định ban đầu. Sửa theo trực giác sẽ dồn công nhầm vào Benign trong khi nút thắt thực nằm ở PortScan.

\textbf{Ngưỡng tin cậy là đòn bẩy rẻ nhất khi cảnh báo giả là dự đoán độ tin cậy thấp.} Phần lớn false positive của Benign là các dự đoán tấn công sát ranh giới quyết định. Một ngưỡng duy nhất ($T=0{,}6$) nâng Benign Recall từ 75,4\% lên 95,2\% mà không phải huấn luyện lại và không làm tụt recall của các lớp tấn công vốn được dự đoán với độ tin cậy cao.

\textbf{Cân bằng lớp quan trọng hơn khối lượng thô.} Phiên BruteForce thu được 200.000 flow nhưng không đóng góp gì hơn một tập con cân bằng khoảng 8.000 flow; thu thừa một lớp chỉ làm lệch corpus và tốn công xử lý.

\end{document}
